\section{VLM 的一般性架构}

\subsection{三大组件}

标准的 VLM 架构包含三个核心组件：

\begin{enumerate}
  \item \textbf{视觉编码器}（Vision Encoder）：将图像/视频编码为视觉 token。常用 ViT (Vision Transformer)
  \item \textbf{投影/对齐模块}（Projector/Adapter）：将视觉 token 映射到语言模型的 embedding 空间
  \item \textbf{语言模型骨干}（LLM Backbone）：处理混合的文本 + 视觉 token 序列
\end{enumerate}

\subsection{动态分辨率处理}

现代 VLM 普遍支持动态分辨率输入，不再将所有图像调整到固定大小：
\begin{itemize}
  \item 将图像分割成多个 patch
  \item 每个 patch 编码为若干视觉 token
  \item 更高分辨率 $\to$ 更多视觉 token $\to$ 更长的序列
\end{itemize}

\subsection{本章小结}

VLM = 视觉编码器 + 投影模块 + 语言模型。动态分辨率处理是当前的标准做法，使模型能够处理不同大小的图像。

